\section{Introduction}\label{sec:introduction}

A retail-sized payment from a sender in the United States to a recipient in Mexico routed through the SWIFT correspondent banking network carries an all-in cost of approximately 330 basis points of the transferred amount, according to the World Bank's remittance prices panel \citep{rpw_q3_2025}. The same payment routed through a USDT transfer on the Tron network carries a per-transaction fee of approximately 37 basis points at a transfer size of USD 400 \citep{tronsave_blog_2026}. The roughly ninefold gap has been the empirical underpinning of the recent argument that stablecoins represent a frontier technology in cross-border retail payments. The comparison is incomplete, however, because it compares the announced cost of the stablecoin alternative against the all-in cost of the incumbent.

The sender of a stablecoin payment does not pay only the on-chain transfer fee. The sender holds, however briefly, a private-issuer liability that is convertible into the underlying fiat unit of account at par on a best-effort basis. The promise of par convertibility can fail in two ways. The issuer may not have enough liquid reserves of sufficient credit quality to honour redemption at par, in which case the holder bears a counterparty loss. The market for converting the liability outside the issuer's primary venue has finite depth, so the realised conversion price falls below par when demand for conversion is high, in which case the holder bears a depeg loss. The depeg loss is itself a coordination outcome: when many holders convert at once, the issuer's primary redemption is exhausted and the secondary price falls, a self-fulfilling run that has to be priced with a coordination model. Counterparty loss is a distinct, issuer-level component; the two are modelled jointly, and at the calibration they price additively, with their interaction second-order. A complementary fragility arises on the settlement rail itself, where blockchain congestion can destabilise even a perfectly reserve-backed digital money \citep{klee_etal_2026_fragility}; we take the rail's transfer cost as given and price the liability the sender briefly carries.

We price the premium in two components. The counterparty component is an issuer-level credit exposure, the product of a default hazard and a loss given default read from market proxies for the issuer's reserve quality; it does not depend on what other holders do. The depeg component is a coordination outcome, which we price with a global game in the spirit of \citet{morris_shin_1998}, adapted to a dual-channel exit through the issuer's primary redemption and a secondary market with depth-dependent price impact; Figure~\ref{fig:global_game} previews the coordination structure. The model is solved at the endogenous fixed point: each holder's redemption-threshold strategy is consistent with the aggregate conversion mass it implies, and the marginal holder's indifference condition integrates over the latent issuer fundamental with the secondary price formed jointly. The depeg component splits, in turn, into a baseline-regime and a stress-regime piece, so the structural cost the holder bears decomposes into three numbers in basis points whose probability-weighted sum is the risk premium, denoted $\Pi_{\text{risk}}$, that we add to the on-chain transfer fee to recover the risk-adjusted cost of the stablecoin payment.\footnote{We use the term risk premium for brevity; the object is an expected-loss adjustment, the probability-weighted loss a risk-neutral holder expects, with no compensation for risk aversion added. This is also why a positive adjustment is consistent with the coin quoting at par in normal times: the expected loss is borne at conversion through the off-ramp and in rare stress states, not as a continuous discount on the quoted price, which arbitrage holds at par until redemption pressure binds.}

We calibrate the model using hourly price data for the two largest fiat-backed stablecoins around their two largest historical depeg episodes: USDC during the March 2023 Silicon Valley Bank failure, and USDT during the May 2022 Terra collapse. The two episodes are treated separately because the two assets differ in features that matter for the calibration: USDC and USDT have different reserve composition, different primary redemption arrangements, and different secondary-market venues. These differences translate into distinct calibration requirements for the two stablecoins while sharing the same coordination-game structure.

The two calibrations produce risk premiums of 49 basis points for USDC and 71 for USDT, and these two numbers carry the paper's three conclusions. First, the premium is material, of the same order as the on-chain transfer fee it sits on top of, yet small enough that it does not by itself overturn the corridor rankings: pricing a stablecoin payment at its announced fee alone understates its cost by a structural margin that no operational improvement removes. Second, the two coins are different assets rather than one category, with premiums inverted in composition, USDC's depeg-dominant (59\% of the total across the two depeg components) and USDT's counterparty-dominant (84\%), because the same coordination machinery applied to two issuers with different reserves and secondary venues returns two distinct risk profiles, so a single supervisory rule would misprice one of them. Third, the advantage the stablecoin enjoys where it wins is contingent on the absence of an upgraded public payment rail rather than a structural property of the technology, as the corridor comparison shows.

We then add the risk premium and the empirical off-ramp basis to the on-chain transfer fee and benchmark the result against publicly operated instant payment systems, correspondent banking, and fintech remittance providers on three corridors that span the policy-relevant space: an intra-Brazilian retail payment of USD 40 equivalent, a US-to-Mexico remittance of USD 400, and an EU-to-Brazil cross-border transfer of USD 1{,}000. The three corridors trace three regimes. Stablecoins are dominated where an upgraded public domestic instant payment system has already been built, as in the Brazilian retail corridor where PIX clears at zero cost. Stablecoins win net of the structural risk premium where no upgraded public cross-border alternative exists, as on the US-Mexico and EU-Brazil corridors today, where the comparable fintech-upgraded alternative is the closest competitor. Finally, an upgraded public cross-border alternative dominates the stablecoin path once one is operational: a BIS Project Nexus rail going live in 2027 at the cost structure of the SEPA Instant Payments Regulation would dominate by margins of one to two orders of magnitude on both cross-border corridors.

Our contribution is to turn that exposure into a price, issuer by issuer. We model the depeg as a global game on the issuer's liability, with a secondary market that is convex and contracts under stress, and we calibrate it separately for USDC and for USDT to recover a structural risk premium split into its counterparty and depeg parts. That decomposition is what the rest of the paper runs on. It tells a depeg-dominated coin from a counterparty-dominated one, gives each a magnitude in basis points that enters the corridor comparison against public, incumbent, and fintech rails, and maps onto the two instruments a regulator would actually use. The framework is not tied to these two coins or three corridors: it applies to any issuer and corridor for which a calibration episode and cost data exist.

The question matters beyond the measurement exercise. Stablecoins are scaling fastest on precisely the cross-border retail corridors where public infrastructure is weakest, and the policy debate has priced them at their announced on-chain fee, the one number that omits the structural cost the holder actually bears. Pricing that cost, and recognising that it differs by issuer, reframes two decisions a central bank faces: whether the right response to stablecoin adoption is to build or to license cross-border payment infrastructure, and whether a single rulebook can supervise instruments whose risks are composed so differently.

The remainder of the paper is organised as follows. Section~\ref{sec:literature} reviews the related literature. Section~\ref{sec:risk_factors} sets out the two risk factors of a stablecoin, the counterparty credit exposure and the depeg coordination problem. Section~\ref{sec:model} presents the global game and its equilibrium properties. Section~\ref{sec:calibration} calibrates the premium, reports its decomposition for the two stablecoins, and checks the model's depeg mechanism against the data. Section~\ref{sec:empirical_comparison} presents the corridor-level comparison against public, incumbent, and fintech rails across three use cases. Section~\ref{sec:discussion} draws the policy implications, both the heterogeneity of the asset class and the public-infrastructure trajectory. Section~\ref{sec:conclusion} concludes. Appendix~\ref{app:proofs} contains the proofs and the numerical implementation; Appendix~\ref{app:calib_details} documents the calibration mechanics and the identification check; Appendix~\ref{app:sensitivity} reports the regulatory-uncertainty band; Appendix~\ref{app:cost_data} documents the payment-rail cost data. A separate internet appendix collects the counterparty-parameter sensitivity, the calibration residuals, descriptive evidence on stablecoin pricing around regulatory enactments, a two-class robustness extension, and a consolidated summary of the corridor comparison.
